\subsection{Deep Reinforcement Learning}

Reinforcement Learning (RL) formalises the decision-making problem as a Markov Decision
Process (MDP), defined by the tuple \((\mathcal{S}, \mathcal{A}, P, R, \gamma)\), where
\(\mathcal{S}\) is the state space, \(\mathcal{A}\) is the action space, \(P\) is the
transition probability, \(R\) is the reward function, and \(\gamma\) is the discount factor.

\subsubsection{Deep Q-Networks (DQN)}

Standard Q-Learning stores state-action values in a table, which is computationally
intractable for continuous or high-dimensional state spaces (such as a UAV's position).
Mnih et al.~\cite{mnih2015humanlevel} introduced the Deep Q-Network (DQN), which solves
this by using a neural network to approximate the Q-value function \(Q(s, a; \theta)\),
where \(\theta\) represents the network weights. The agent minimises the temporal difference
error defined by the Bellman equation:
\begin{equation}
    L(\theta) = \mathbb{E}\left[\left(r + \gamma \max_{a'} Q(s', a'; \theta^{-}) - Q(s, a; \theta)\right)^2\right] \label{eq:dqn_loss}
\end{equation}
This project implements the DQN agent via the Stable-Baselines3
library~\cite{raffin2021stable} within a custom Gymnasium environment~\cite{towers2024gymnasium},
learning the non-linear mapping between the complex network state (RSSI, AoI, battery) and
the optimal flight vector.
